\ifdefined\zadaniaJednaNacitane \endinput \fi
\def\zadaniaJednaNacitane{1}

%%%%%%%%%%%%%%%%%%%%%%%%%%%%%%%%%%%%%%%%%%%%%%%%%%%%%%%%%%%%%%%%%%%%%%%%%%%%%%%%
%prvá séria
%%%%%%%%%%%%%%%%%%%%%%%%%%%%%%%%%%%%%%%%%%%%%%%%%%%%%%%%%%%%%%%%%%%%%%%%%%%%%%%%

\newcommand{\zadIsI}{
V trojuholníku \(ABC\) os uhla \(ABC\) pretína \(AC\) v bode \(D\). Rovnobežka s \(BC\) vedená bodom \(D\) pretína \(AB\) v bode \(E\). Platí, že strany \(AB\) a \(BC\) sú v pomere \(2:1\). V akom pomere delí \(E\) úsečku \(AB\)?
}

%%%%%%%%%%%%%%%%%%%%%%%%%%%%%%%%%%%%%%%%%%%%%%%%%%%%%%%%%%%%%%%%%%%%%%%%%%%%%%%%

\newcommand{\zadIIsI}{
Dokážte, že súčin troch po sebe idúcich kladných celých čísel nemôže byť treťou mocninou kladného celého čísla.
}

%%%%%%%%%%%%%%%%%%%%%%%%%%%%%%%%%%%%%%%%%%%%%%%%%%%%%%%%%%%%%%%%%%%%%%%%%%%%%%%%

\newcommand{\zadIIIsI}{
\textit{Obec} je súvislý útvar zložený zo štvorcov rovnakej veľkosti, ktoré navzájom susedia hranou. Na obrázku pod úlohou je uvedené, ako by napríklad mohla vyzerať \textit{obec} zložená z ôsmich štvorcov.

\textit{Mestom} nazveme takú \textit{obec}, ktorá sa skladá z aspoň \(2026\) štvorcov. \textit{Nedeliteľné mesto} je také, ktoré sa nedá rozdeliť na dve alebo viacero \textit{miest}. Aký je maximálny počet štvorcov, z ktorých sa môže skladať \textit{nedeliteľné mesto}?

\begin{center}
    \includegraphics[scale=0.6]{obr/4-3.PNG}
\end{center}

}

%%%%%%%%%%%%%%%%%%%%%%%%%%%%%%%%%%%%%%%%%%%%%%%%%%%%%%%%%%%%%%%%%%%%%%%%%%%%%%%%

\newcommand{\zadIVsI}{
Pre aké najmenšie celé číslo \(n>1\) existuje kladné celé číslo \(m\) také, že číslo \(m^{1000}+n^{1000}+n\) je deliteľné číslom \(mn\)?
}

%%%%%%%%%%%%%%%%%%%%%%%%%%%%%%%%%%%%%%%%%%%%%%%%%%%%%%%%%%%%%%%%%%%%%%%%%%%%%%%%

\newcommand{\zadVsI}{
Na tabuli je napísaných \(n\) núl v jednom riadku. Hráme hru. V každom ťahu zotrieme jednu z cifier na tabuli a nahradíme ju inou cifrou na tom istom mieste. Dodržiavame pritom tieto obmedzenia:
\begin{itemize}
    \item Píšeme iba cifry \(0\), \(1\) a \(2\).
    \item Rozdiel napísanej a zotretej cifry musí byť \(1\).
    \item Cifru nemožno zotrieť, ak sa niekde napravo od nej nachádza rovnaká cifra.
    \item Nemožno napísať takú cifru, ktorá sa už nachádza niekde napravo od práve zotretej cifry.
\end{itemize}
Určte, na koľko najmenej ťahov môžeme dosiahnuť stav, že na tabuli budú napísané samé jednotky.
}

%%%%%%%%%%%%%%%%%%%%%%%%%%%%%%%%%%%%%%%%%%%%%%%%%%%%%%%%%%%%%%%%%%%%%%%%%%%%%%%%

\newcommand{\zadVIsI}{
Nájdite všetky funkcie \(f: \mathbb{Z}\rightarrow\mathbb{Z}\) také, že pre všetky \(a,b\in\mathbb{Z}\) platí:
\[f(a+b)+f(ab-1)=f(a)f(b)+2.\]
}

%%%%%%%%%%%%%%%%%%%%%%%%%%%%%%%%%%%%%%%%%%%%%%%%%%%%%%%%%%%%%%%%%%%%%%%%%%%%%%%%
